\documentclass[11pt,a4paper]{article}

\usepackage{kotex}
\usepackage{fontspec}
\setmainfont{Times New Roman}
\setmainhangulfont{AppleMyungjo}
\setsanshangulfont{Apple SD Gothic Neo}

\usepackage[a4paper,top=18mm,bottom=18mm,left=20mm,right=20mm]{geometry}
\usepackage{fancyhdr}
\setlength{\headheight}{24pt}
\usepackage{enumitem}
\usepackage{mdframed}
\usepackage{array}
\usepackage{booktabs}

\setlength{\parindent}{1em}
\linespread{1.18}

\pagestyle{fancy}
\fancyhf{}
\renewcommand{\headrulewidth}{0.6pt}
\fancyhead[L]{\sffamily\bfseries 2026 고1 영어 변형 --- 지문 35}
\fancyhead[R]{\sffamily [디지털 미니멀리즘]}
\renewcommand{\footrulewidth}{0pt}
\fancyfoot[C]{\framebox[16mm][c]{\thepage}}

\newcommand{\qno}[1]{\par\vspace{8pt}\noindent{\sffamily\bfseries #1.}\ }
\newcommand{\instr}[1]{{\sffamily #1}\par\vspace{3pt}}
\newlist{choices}{enumerate}{1}
\setlist[choices]{label=,leftmargin=1.5em,itemsep=2pt,topsep=3pt,parsep=0pt}
\newmdenv[linewidth=0.6pt,roundcorner=3pt,innertopmargin=7pt,innerbottommargin=7pt,
          innerleftmargin=9pt,innerrightmargin=9pt,skipabove=5pt,skipbelow=5pt]{pbox}
\newcommand{\nemo}[1]{\fbox{\,#1\,}}

\begin{document}

\begin{center}
{\fontsize{15}{17}\selectfont\sffamily\bfseries [지문 35] 다음 글을 읽고 물음에 답하시오.}
\end{center}
\vspace{4pt}

\begin{pbox}
It is evident through various case studies that adopting digital minimalism can lead to
improved \rule{2cm}{0.4pt}. ( ① ) For instance, a study conducted by researchers at
Stanford University found that participants who limited their use of social media and digital
devices reported lower levels of stress, anxiety, and depression. ( ② ) By reducing the
constant stream of information and notifications, individuals were able to focus more on
meaningful interactions and activities that promoted well-being. ( ③ ) Additionally,
participants reported relying on social media to get information for their purchasing
decisions. ( ④ ) In another case, a corporate professional struggling with burnout and
insomnia implemented digital minimalism techniques, such as setting boundaries for phone
usage and practicing mindfulness. ( ⑤ ) As a result, he experienced improved sleep quality,
reduced stress levels, and a newfound sense of calm and clarity.

\vspace{4pt}
\noindent\footnotesize{* insomnia: 불면증}
\end{pbox}

\bigskip

\qno{1}\instr{다음 글에서 전체 흐름과 관계 없는 문장은?}
\begin{choices}
\item ① \quad ② \quad ③ \quad ④ \quad ⑤
\end{choices}

\bigskip

\qno{2}\instr{다음 빈칸에 들어갈 말로 가장 적절한 것은?}
\begin{choices}
\item ① productivity at work
\item ② digital literacy skills
\item ③ consumer decision-making
\item ④ psychological well-being
\item ⑤ social media engagement
\end{choices}

\bigskip

\qno{3}\instr{윗글의 주제로 가장 적절한 것은?}
\begin{choices}
\item ① the growing popularity of digital minimalism among young consumers
\item ② evidence that reducing digital usage can benefit mental health
\item ③ how social media affects purchasing behavior in the digital age
\item ④ techniques for practicing mindfulness in a connected world
\item ⑤ the relationship between screen time and academic performance
\end{choices}

% ===================== Q4 =====================
\qno{4}\instr{[서술형] 윗글의 \textit{setting boundaries for phone usage}가 의미하는 바를 본문 맥락에 맞게 우리말로 쓰시오.}
\par\vspace{4pt}\noindent
$\Rightarrow$ \rule{\linewidth}{0.4pt}

\end{document}
